\begin{abstract}
Encrypted-traffic detectors publish accuracy with no stated limit, so a
practitioner cannot say what pairwise error is impossible for a given
observation window. We measure a two-sided $\epsilon$-indistinguishability
bracket on real QUIC wire data, one certificate per service pair, whose upper
side is a finite-sample bound on the total variation distance between that
pair's prefix laws and whose lower side is a classifier-based bound at a nominal
confidence level. The certified object is pairwise: its decision-theoretic
content is the two-point bound on a test between those services, and we claim no
multi-class or open-set reduction.

At the coarsest pre-registered channel, 336 of 990 service pairs (33.9\%) admit
such a certificate, the upper side being the scarce one at the rung where the
gate is defined; because the pairs form a complete graph over 45 services they
are not independent, and a service-cluster interval gives $[26.5\%, 43.1\%]$.
Raising the alphabet resolution is associated with higher divergence and
destroys the certificate at the same time, since the confidence radius grows with
the alphabet while the threshold does not; the upper-side set falls from 340 pairs
to none four rungs finer, and the boundary has a closed form. Where the
certificate is lost, a divisibility bound keeps the large-prefix regime
non-empty under a realizability assumption. On an independent corpus the lower
side reproduces almost exactly (0.9967 against 0.9960) while the joint gate
falls to 15.75\% under that corpus's capped lower-side convention; the gap is
significant before matching but not statistically resolved at matched evidence.

We validate the bracket against real detectors: on the 340 pairs in the
informative stratum, the observed floor-violation rates of a three-model panel
(0.6--2.9\%) fall below the 5\% null rate ($p = 0.9765$ at the strongest cell;
under the $2.5\%$ null implied by the $0.975$ joint level, $p = 0.346$), so no
detector's closed-set violation rate exceeds the certified floor. The compared
quantity is a closed-set error surrogate and we scope the claim to it. The paper also reports one methods and
limits finding: the standard chain-rule estimator returns sublinear growth on a
source whose divergence is known, so it cannot adjudicate whether divergence
grows sublinearly across dependent packet positions.
\end{abstract}
